% !TEX root = ../main.tex
\section{Methodology}

\subsection{Reliability Metrics}

We define LLM reliability as a multi-dimensional construct captured by four key metrics:

\begin{itemize}
    \item \textbf{Success Rate} ($S$): The proportion of evaluations where the model produces a correct or acceptable output, defined as $S = \frac{1}{N}\sum_{i=1}^{N} \mathbb{1}[\text{success}_i]$, where $N$ is the total number of evaluations for a given agent-benchmark pair.

    \item \textbf{Perturbation Robustness} ($P$): The model's ability to maintain correct outputs when inputs are perturbed, computed as $P = \frac{1}{K}\sum_{j=1}^{K} S_j$, where $S_j$ is the success rate under perturbation type $j$, and $K$ is the number of perturbation types.

    \item \textbf{Fault Tolerance} ($F$): The model's ability to recover from injected faults, measured as $F = \frac{1}{M}\sum_{k=1}^{M} R_k$, where $R_k \in \{0,1\}$ indicates successful recovery from fault $k$, and $M$ is the number of fault injections.

    \item \textbf{Repeated-Run Consistency} ($C$): The consistency of model outputs across multiple runs with identical inputs, measured using pairwise agreement: $C = \frac{2}{N(N-1)}\sum_{i < j} \mathbb{1}[\text{outcome}_i = \text{outcome}_j]$, where $N$ is the number of repeated runs.
\end{itemize}

The \textbf{Composite Reliability Score} ($R$) is defined as the weighted combination:
\[
R = \alpha S + \beta P + \gamma F + \delta C
\]
where $\alpha = \beta = \gamma = \delta = 0.25$ by default, though alternative weightings can be explored.

\subsection{Fault Injection Protocol}

We inject six types of faults during evaluation to simulate production-environment anomalies:

\begin{enumerate}
    \item \textbf{Artificial Timeout}: Simulates a model response that exceeds a latency threshold, requiring retry logic.
    \item \textbf{Temporary API Failure}: Simulates transient API unavailability, testing retry mechanisms.
    \item \textbf{Invalid Model Response}: Injects a malformed response to test output validation.
    \item \textbf{Tool Failure}: Simulates a tool execution error (e.g., web search unavailability).
    \item \textbf{Network Interruption}: Simulates a connection reset during API communication.
    \item \textbf{Context Truncation}: Simulates prompt truncation due to context window limits.
\end{enumerate}

Each fault type is injected systematically across all model-benchmark combinations, with recovery outcomes recorded as success (model produces valid output after recovery) or failure (model produces incorrect output or errors out).

\subsection{Perturbation Types}

We apply five perturbation types to evaluate input robustness:

\begin{enumerate}
    \item \textbf{Whitespace}: Addition or modification of whitespace characters in prompts.
    \item \textbf{Synonym Substitution}: Replacement of words with semantically equivalent alternatives.
    \item \textbf{Reordering}: Rearrangement of sentence or instruction order.
    \item \textbf{Prompt Wrapper}: Addition or modification of instruction wrappers.
    \item \textbf{Formatting}: Changes to output format specifications.
\end{enumerate}

\subsection{Ranking Methodology}

We employ two ranking approaches:

\paragraph{Per-Criterion Ranking:} For each reliability dimension (success, reliability, composite), models are ranked within each benchmark, producing ordered lists from best (rank 1) to worst.

\paragraph{Borda Count Aggregation:} To produce a consensus ranking, we apply the Borda count method. For each ranking criterion, a model receives $B_i = N_{\text{models}} - \text{rank}_i$ points. The total Borda score is summed across all criteria and normalized to $[0, 1]$.

\subsection{Experimental Setup}

Our evaluation covers 6 API-based agents and 15 additional agents across 4 benchmarks, with 3--8 repeated runs per agent-benchmark pair (API-based models: 3 runs; locally hosted models: 7--8 runs). In total, we conducted 142 experiments totaling 31,981 individual evaluations.
